\chapter{Background}
\label{ch:bg}

\section{Learning With Errors}

Given a public matrix $A$ and a vector $b = As + e \pmod q$, where $s$ is
secret and $e$ is a vector of small random ``errors'', the Learning With Errors
(LWE) problem~\cite{regev} asks for $s$. Without $e$ this is linear algebra.
With it, the best known classical and quantum algorithms take exponential
time, and Regev showed that solving random instances is as hard as
worst-case lattice problems. Module-LWE~\cite{langlois} replaces the entries
of $A$ with polynomials in $R_q = \mathbb{Z}_q[X]/(X^{256}+1)$, which shrinks
keys and allows fast multiplication with the number-theoretic transform. Both
ML-KEM and ML-DSA use $n = 256$; ML-KEM uses $q = 3329$ and ML-DSA uses
$q = 8\,380\,417$.

\section{ML-KEM (FIPS 203)}

A key-encapsulation mechanism has three algorithms:
\begin{align*}
(ek, dk) &\leftarrow \mathsf{KeyGen}() \\
(K, c) &\leftarrow \mathsf{Encaps}(ek) \\
K &\leftarrow \mathsf{Decaps}(dk, c)
\end{align*}
\textsf{Encaps} chooses a random 32-byte message $m$, derives $(K, r) =
G(m \parallel H(ek))$, and encrypts $m$ under $ek$ with randomness $r$.
\textsf{Decaps} decrypts to $m'$, re-derives $(K', r')$, re-encrypts and
compares. If the ciphertexts differ it returns a pseudorandom key derived from
a secret value $z$ instead of an error (implicit rejection), so an attacker
learns nothing from malformed ciphertexts.

Two properties matter for protocol design. First, \textbf{a KEM is not
Diffie--Hellman}: only the receiver has a key pair, and the sender's
contribution is fresh randomness. A KEM therefore authenticates neither party
by itself; authentication must come from signatures. Second, \textbf{a KEM
has no signing operation}, which is the root of defect D1 in the base paper.

\begin{table}[h]
\centering
\caption{ML-KEM parameter sets (FIPS 203)}
\label{tab:mlkem}
\begin{tabular}{lrrrrr}
\toprule
Parameter set & $k$ & category & $ek$ (B) & $dk$ (B) & ciphertext (B) \\
\midrule
ML-KEM-512 & 2 & 1 & 800 & 1,632 & 768 \\
ML-KEM-768 & 3 & 3 & 1,184 & 2,400 & 1,088 \\
ML-KEM-1024 & 4 & 5 & 1,568 & 3,168 & 1,568 \\
\bottomrule
\end{tabular}
\end{table}

The encapsulation key is $k$ polynomials of 256 coefficients at 12 bits each,
plus a 32-byte seed for $A$: for ML-KEM-1024, $4 \times 384 + 32 = 1{,}568$
bytes. The shared secret is always 32 bytes.

\section{ML-DSA (FIPS 204)}

ML-DSA signs with a Fiat--Shamir-with-aborts construction. The signer picks a
masking vector $y$, computes $w = Ay$, derives a challenge $c$ from the hash of
the message representative $\mu$ and the high bits of $w$, and outputs $z = y
+ cs_1$ only if $z$ leaks nothing about the secret $s_1$; otherwise it starts
again. The verifier recomputes the high bits of $Az - ct_1$ from the public
key, hashes them with $\mu$, and checks the result against $c$. The message is
never inside the signature: the verifier needs the message, the signature and
the public key. FIPS~204 adds a \emph{context string} to the signed message;
we use \code{qrcav-v1} for domain separation.

\begin{table}[h]
\centering
\caption{ML-DSA parameter sets (FIPS 204) and ECDSA P-256 for comparison}
\label{tab:mldsa}
\begin{tabular}{lrrr}
\toprule
Scheme & category & public key (B) & signature (B) \\
\midrule
ML-DSA-44 & 2 & 1,312 & 2,420 \\
ML-DSA-65 & 3 & 1,952 & 3,309 \\
ML-DSA-87 & 5 & 2,592 & 4,627 \\
ECDSA P-256 & -- & 65 & 70--72 \\
\bottomrule
\end{tabular}
\end{table}

\section{Symmetric building blocks}

\paragraph{Hash.} $H$ is SHAKE-256 over length-prefixed inputs, so
$H(ab, c) \neq H(a, bc)$.

\paragraph{Key derivation.} HKDF~\cite{hkdf} with HMAC-SHA3-256 extracts a
pseudorandom key from the KEM shared secret and expands it into independent
keys and IVs for each direction, each with its own label.

\paragraph{Records.} AES-256-GCM~\cite{gcm} provides confidentiality and
integrity. Its one hard rule is that a (key, nonce) pair must never repeat:
reuse reveals the authentication key and allows forgeries. We derive the
nonce as $v \oplus (0^{32} \parallel \mathrm{be64}(ctr))$ from a per-direction
IV $v$ and a record counter.

\section{PBFT}

PBFT~\cite{pbft} runs on $n \ge 3f + 1$ replicas, one of which is the leader
(primary) for the current view. The leader assigns a sequence number to a
request in a PRE-PREPARE; replicas broadcast PREPARE, and after $2f+1$
matching prepares broadcast COMMIT; after $2f+1$ commits they execute and
reply. Any two quorums of $2f+1$ intersect in at least one honest replica,
which gives \emph{safety} (no two honest replicas commit different values at
one sequence number). \emph{Liveness} needs the leader to make progress; if it
does not, replicas time out and run a view change that installs the next
leader and carries prepared requests across. Messages per request grow as
$O(n^2)$.

\section{IEEE 802.11p}

IEEE~802.11p~\cite{ieee80211p} is Wi-Fi adapted for vehicles: 10\,MHz
channels at 5.9\,GHz, typically 6\,Mb/s for safety messages, no association
(communication outside the context of a BSS), and EDCA access categories that
let safety traffic (AC\_VO) win the channel before bulk traffic (AC\_BE).
Periodic beacons are link-layer broadcasts with no acknowledgement or
retransmission. A 1,400-byte payload occupies the channel for about 2\,ms at
6\,Mb/s, so a 20\,KB handshake costs about 30\,ms of airtime before any
contention or loss.
